
\chapter{Zakończenie}

Celem pracy było zbadanie skuteczności rozwiązań sztucznej inteligencji
w~projektowaniu płytek PCB, na przykładzie trasowania połączeń
i~rozmieszczania elementów. Przeprowadzone eksperymenty pokazują, że
projektowanie złożonych płytek drukowanych pozostanie na razie w~rękach
inżynierów. W~przypadku trasowania wkrótce może jednak nadejść moment,
w~którym wspieranie się sztuczną inteligencją przy wyznaczaniu ścieżek
stanie się codziennością. Sieć neuronowa już teraz wyraźnie przyspiesza
wyszukiwanie ścieżek, choć ich jakość nie dorównuje jeszcze klasycznemu
autorouterowi. W~przypadku rozmieszczania elementów upłynie natomiast
znacznie więcej czasu, zanim doświadczenie i~wiedza projektanta ustąpią
miejsca sztucznej inteligencji.

Nie oznacza to, że temat powinien przestać interesować środowisko
naukowe. Praca stanowi solidną podstawę do dalszych badań, a~kierunków
ich rozwoju jest wiele. Można nawet podjąć próbę wykorzystania dużych
modeli językowych, ponieważ projekty płytek zapisywane są w~formatach
tekstowych, takich jak pliki KiCad czy Specctra DSN. Na dzień dzisiejszy
pytanie o~skuteczność sztucznej inteligencji w~projektowaniu płytek PCB
pozostaje więc otwarte.
